\scanpage{43}
où
\[
\unsure{\Theta|\theta}:\underline O_X^*\longrightarrow\underline{\Phi}_X\mathbin{\tikz[baseline=(op.base)]{\node[inner sep=0pt] (op) {$\times$};\draw[overlay,thin] (-.33,-.22) .. controls (.09,-.33) and (-.04,.48) .. (.35,.49);}}R_X^*
\]
est défini par
\[
\Theta(f)=(-\varphi_f\mathpunct{\tikz[baseline=(op.base)]{\node[inner sep=0pt] (op) {$,$};\draw[overlay,thin] (-.33,-.22) .. controls (.09,-.33) and (-.04,.48) .. (.35,.49);}}f^{\deleted{\unsure{-1|1}}})
\]
On a donc deux suites exactes
\diagram{\begin{tikzcd}[ampersand replacement=\&,column sep=large]
\underline{\Phi}_X\arrow[r] \& \underline{\mathrm{QF}}_X\arrow[r,"{\mathrm{div}}"] \& \overset{\simeq\,\underline{R}_X^*/\underline{O}_X^*}{\underline{\mathrm{Div}}_X}\arrow[r]\arrow[rrd,dash,bend right=12] \& 0 \& \\ \& \& \& \& \text{diviseurs de \uline{Cartier}}
\end{tikzcd}}
\diagram{\begin{tikzcd}[ampersand replacement=\&,column sep=large]
\underline R_X^*\arrow[r,"{f\mapsto\unsure{\zeta_f|\xi_f|\{f\}}}"] \& \underline{\mathrm{QF}}_X\arrow[r] \& \underline{\Phi}_X/\operatorname{Im}\underline{O}_X^*\arrow[r] \& 0
\end{tikzcd}}

Les \inserted{sections sur $X$ du} noyau de $\underline{\Phi}_X\unsure{-|\to}\underline{\mathrm{QF}}_X$ est formé des fonctions admissibles $\varphi$ telles qu'il existe $f\in\Gamma(X,\underline O_X^*)$ avec $(\deleted{\unsure{0}}1,\varphi)=(f,-\varphi_f)$, ce qui exige $f=1$ et $-\varphi_f(\unsure{(P)v|(P)f|P,v})=v(f(P))=0$ d'où $\varphi=0$, donc $\underline{\Phi}_X-\underline{\mathrm{QF}}_X$ est injectif::

\uline{les quasi-fonctions de diviseur nul sont les fonctions admissibles sur} $X$.

Les \inserted{sections sur $X$ du} noyau de $\underline R_X^*-\underline{\mathrm{QF}}_X$ est formé des $f\in{\inserted{\Gamma R_X^*}}$ telles qu'il existe $g\in\Gamma(X,\underline O_X^{\unsure{\times|*}})$, t.q. $(f,0)=(g,-\varphi_g)$. On doit avoir
% archive sheet number: 42
